\documentclass[11pt]{article}
\usepackage[margin=0.85in]{geometry}
\usepackage{amsmath,amssymb,amsthm}
\usepackage{graphicx}
\usepackage{microtype}
\usepackage[colorlinks=true,linkcolor=blue,urlcolor=blue,citecolor=blue]{hyperref}

\newtheorem{theorem}{Theorem}
\newtheorem{remark}{Remark}
\newcommand{\R}{\mathbb R}
\newcommand{\domn}{\operatorname{dom}}
\newcommand{\ic}{\mathbin{\square}}

\title{A disk-lens counterexample to Lipschitz regularity\\
under infimal convolution}
\author{Candidate solution packet for arXiv:1710.08233}
\date{17 August 2026}

\begin{document}
\maketitle

\paragraph{Status.}
Candidate counterexample, likely valid, pending human review.

\paragraph{Source.}
Simon Zugmeyer, \emph{Sharp trace Gagliardo--Nirenberg--Sobolev inequalities
for convex cones, and convex domains}, arXiv:1710.08233 \cite{Zugmeyer},
Section~2.2, source PDF page~10.  For extended-real functions, the source uses
\[
 (f\ic g)(x)=\inf_{y\in\R^n}\{f(y)+g(x-y)\},\qquad
 \domn f=\{x:f(x)<+\infty\}.
\]
After a nonconvex, empty-interior example, it conjectures that Lipschitz
functions on convex domain(s) with nonempty interior have a Lipschitz infimal
convolution.

\begin{center}
\includegraphics[width=0.96\textwidth]{figures/open_problem_crop.png}
\end{center}

\section*{Proof intuition}

For two unit disks, the feasible decompositions of a point $x$ form the lens
$D\cap(x-D)$.  Near the boundary of $D+D$, this lens closes tangentially:
its vertical radius is proportional to the square root of the distance to the
boundary.  An affine function that reads the vertical coordinate makes the
infimal convolution select the bottom of the lens.  The resulting value
function therefore has a square-root cusp, even though both input functions
are Lipschitz on the same smooth compact convex domain.

\section*{Counterexample}

Let
\[
 D=\{(u_1,u_2)\in\R^2:u_1^2+u_2^2\leq 1\}
\]
and define $f,g:\R^2\to\R\cup\{+\infty\}$ by
\[
 f(u_1,u_2)=
 \begin{cases}1+u_2,&(u_1,u_2)\in D,\\+\infty,&\text{otherwise},\end{cases}
 \qquad
 g(u_1,u_2)=
 \begin{cases}0,&(u_1,u_2)\in D,\\+\infty,&\text{otherwise}.
 \end{cases}
\]

\begin{theorem}
The functions $f$ and $g$ have the same compact convex domain $D$, whose
interior is nonempty.  Their restrictions to $D$ are respectively
$1$-Lipschitz and $0$-Lipschitz.  Both extended-real functions are proper,
lower semicontinuous, convex, and nonnegative, and $f\ic g$ is exact.
Nevertheless, $f\ic g$ is not Lipschitz on its domain $D+D=2D$; in fact it is
not locally Lipschitz at $(2,0)$ relative to $2D$.
\end{theorem}

\begin{proof}
All input claims are immediate: $f$ is an affine function plus the convex
indicator of the closed disk, while $g$ is that indicator.  Compactness of
the feasible set gives exactness at every point of $2D$.

For $0\leq t\leq2$, write $x_t=(t,0)$.  A decomposition
$x_t=a+b$ with $a,b\in D$ is equivalent, on writing $a=(s,r)$, to
\begin{equation}\label{eq:constraints}
 s^2+r^2\leq1,
 \qquad
 (t-s)^2+r^2\leq1.
\end{equation}
Thus
\begin{align*}
 (f\ic g)(x_t)
 &=1+\min\{r:(s,r)\text{ satisfies \eqref{eq:constraints}}\}\\
 &=1-\sqrt{1-\min_{s\in\R}\max\{s^2,(t-s)^2\}}.
\end{align*}
The scalar minimax is exact.  Indeed,
\[
 t\leq |s|+|t-s|\leq2\max\{|s|,|t-s|\},
\]
so $\max\{s^2,(t-s)^2\}\geq t^2/4$, with equality for $s=t/2$.
Consequently
\begin{equation}\label{eq:formula}
 (f\ic g)(x_t)=1-\sqrt{1-\frac{t^2}{4}}
 \qquad(0\leq t\leq2).
\end{equation}

For $0<\varepsilon<2$, formula \eqref{eq:formula} gives
\[
 \frac{|(f\ic g)(x_2)-(f\ic g)(x_{2-\varepsilon})|}
      {\|x_2-x_{2-\varepsilon}\|_2}
 =\frac{\sqrt{\varepsilon-\varepsilon^2/4}}{\varepsilon}
 =\sqrt{\frac1\varepsilon-\frac14}\longrightarrow+\infty.
\]
No finite Lipschitz constant exists, even in any relative neighborhood of
$(2,0)$ in $2D$.
\end{proof}

\section*{Wording variants and scope}

\begin{remark}[Open domains]
If ``domain'' was intended to mean an open convex set, replace $D$ by its
open unit disk.  Formula \eqref{eq:formula} holds for $0\leq t<2$.  Both
$x_{2-\varepsilon}$ and $x_{2-2\varepsilon}$ are interior points, while
\[
 \frac{|(f\ic g)(x_{2-\varepsilon})-(f\ic g)(x_{2-2\varepsilon})|}
 {\varepsilon}
 =\frac{\sqrt{2\varepsilon-\varepsilon^2}
 -\sqrt{\varepsilon-\varepsilon^2/4}}{\varepsilon}
 \sim\frac{\sqrt2-1}{\sqrt\varepsilon}.
\]
Hence the convolution is still not globally Lipschitz on its open domain.
\end{remark}

The example refutes the printed global-Lipschitz conjecture.  It does not
refute the weaker assertion that the convolution is locally Lipschitz at each
point of $\operatorname{int}(D+D)$; the closed-disk failure occurs at the
boundary, and the open-disk function is locally Lipschitz at each fixed point
on the displayed diameter.  If a factor is finite and globally $L$-Lipschitz
on all of $\R^n$, then its infimal convolution with any function is
$L$-Lipschitz whenever finite, so the extended-domain interpretation is the
only nontrivial one in the source context.

\section*{Verification and novelty}

The proof is exact and uses no computation.  The bundled checker tests the
scalar minimax inequality with rational arithmetic on 151,151 grid pairs and
prints the divergent quotients for closed and open disks.  All tests pass.

A bounded search on 17 August 2026 covered the four lightweight run indexes,
the exact arXiv id and title, the exact conjecture sentence, the author name,
and close arXiv/web queries combining ``infimal convolution'', ``convex
domain'', ``nonempty interior'', and ``Lipschitz''.  The exact-phrase hit was
the published version of the source itself.  No later answer or matching disk
counterexample was found.  Novelty confidence is moderate: an elementary
example may exist as folklore or under different terminology.

\paragraph{Human-review recommendation.}
Check that the source's ``Lipschitz continuous'' is read globally on the
essential domain, as its wording and surrounding extended-real definitions
indicate.  The formal construction and lens computation are otherwise
elementary.

\begin{thebibliography}{9}
\bibitem{Zugmeyer}
S. Zugmeyer,
\emph{Sharp trace Gagliardo--Nirenberg--Sobolev inequalities for convex
cones, and convex domains},
Ann. Inst. H. Poincar\'e Anal. Non Lin\'eaire 36 (2019), no.~3, 861--885;
arXiv:1710.08233.
\end{thebibliography}

\end{document}
